\documentclass[12pt]{article}

\usepackage{fontspec}
% надёжный шрифт с кириллицей (работает в Overleaf)
\setmainfont{Libertinus Serif}[Ligatures=TeX]
\setsansfont{Libertinus Sans}[Ligatures=TeX]
\newfontfamily\cyrillicfont{Libertinus Serif}
\newfontfamily\cyrillicfontsf{Libertinus Sans}
\usepackage{polyglossia}
\setmainlanguage{russian}
\setotherlanguage{english}

\usepackage{amsmath}
\usepackage{amsfonts}
\usepackage{amssymb}
\usepackage{graphicx}
\usepackage{microtype}

\title{Полный вывод формул оптимального импульсного управления \\ для transfers между соседними орбитами}
\author{На основе работы T. N. Edelbaum}
\date{}

\begin{document}

\maketitle

\section{Фундаментальные уравнения движения}

\subsection{Уравнения в оскулирующих элементах}

Начнем с уравнений изменения элементов орбиты под действием малых возмущений. Для маневров вблизи круговой орбиты используем линеаризованные уравнения.

Введем переменные состояния:
\begin{align*}
x_1 &= \frac{da}{a_0} \quad \text{(безразмерное изменение большой полуоси)} \\
x_2 &= e_y \quad \text{(y-компонента эксцентриситета)} \\
x_3 &= e_x \quad \text{(x-компонента эксцентриситета)} \\
x_4 &= i_y \quad \text{(y-компонента наклонения)} \\
x_5 &= i_x \quad \text{(x-компонента наклонения)}
\end{align*}

Уравнения движения получаются из уравнений Гаусса для оскулирующих элементов:

\begin{align}
\frac{da}{dt} &= \frac{2a^2}{h} \left(e \sin f \cdot F_R + \frac{p}{r} \cdot F_T\right) \\
\frac{de}{dt} &= \frac{1}{h} \left[p \sin f \cdot F_R + ((p+r)\cos f + re) \cdot F_T\right] \\
\frac{di}{dt} &= \frac{r \cos u}{h} F_N \\
\frac{d\Omega}{dt} &= \frac{r \sin u}{h \sin i} F_N
\end{align}

где $F_R$, $F_T$, $F_N$ - компоненты ускорения в радиальной, трансверсальной и нормальной направлениях.

\subsection{Линеаризация около круговой орбиты}

Для круговой орбиты ($e \ll 1$, $r \approx a_0$, $h = \sqrt{\mu a_0}$, $p \approx a_0$):

\begin{align*}
\frac{da}{dt} &\approx 2\sqrt{\frac{a_0}{\mu}} F_T \\
\frac{de_x}{dt} &\approx \sqrt{\frac{a_0}{\mu}} (2F_T \cos \theta + F_R \sin \theta) \\
\frac{de_y}{dt} &\approx \sqrt{\frac{a_0}{\mu}} (2F_T \sin \theta - F_R \cos \theta) \\
\frac{di_x}{dt} &\approx \sqrt{\frac{a_0}{\mu}} F_N \cos \theta \\
\frac{di_y}{dt} &\approx \sqrt{\frac{a_0}{\mu}} F_N \sin \theta
\end{align*}

Вводим новую независимую переменную - интеграл от ускорения тяги:
\begin{equation}
du = |\mathbf{F}| dt
\end{equation}

Тогда компоненты ускорения:
\begin{equation}
F_R = l_R |\mathbf{F}| = l_R \frac{du}{dt}, \quad 
F_T = l_T \frac{du}{dt}, \quad 
F_N = l_N \frac{du}{dt}
\end{equation}

Подставляя, получаем уравнения (1)-(5):

\begin{align}
\frac{dx_1}{du} &= 2\sqrt{\frac{a_0}{K}} l_T \tag{1} \\
\frac{dx_2}{du} &= \sqrt{\frac{a_0}{K}} (2l_T \sin \theta - l_R \cos \theta) \tag{2} \\
\frac{dx_3}{du} &= \sqrt{\frac{a_0}{K}} (2l_T \cos \theta + l_R \sin \theta) \tag{3} \\
\frac{dx_4}{du} &= \sqrt{\frac{a_0}{K}} l_N \sin \theta \tag{4} \\
\frac{dx_5}{du} &= \sqrt{\frac{a_0}{K}} l_N \cos \theta \tag{5}
\end{align}

где $K = \mu$ - гравитационный параметр.

\section{Формулировка вариационной задачи}

\subsection{Гамильтониан и сопряженные переменные}

Гамильтониан системы:
\begin{equation}
H = \sum_{i=1}^5 \lambda_i \frac{dx_i}{du} \tag{6}
\end{equation}

Подставляя уравнения (1)-(5):

\begin{align*}
H = \sqrt{\frac{a_0}{K}} \left[ 
2\lambda_1 l_T + 
\lambda_2 (2l_T \sin \theta - l_R \cos \theta) + 
\lambda_3 (2l_T \cos \theta + l_R \sin \theta) + 
\lambda_4 l_N \sin \theta + 
\lambda_5 l_N \cos \theta
\right]
\end{align*}

Группируем члены с $l_R$, $l_T$, $l_N$:

\begin{align*}
H = \sqrt{\frac{a_0}{K}} \left\{
l_R [-\lambda_2 \cos \theta + \lambda_3 \sin \theta] +
l_T [2\lambda_1 + 2\lambda_2 \sin \theta + 2\lambda_3 \cos \theta] +
l_N [\lambda_4 \sin \theta + \lambda_5 \cos \theta]
\right\}
\end{align*}

\subsection{Вектор праймера}

Вводим компоненты вектора праймера:

\begin{align}
\lambda &= \sqrt{\frac{a_0}{K}} (-\lambda_2 \cos \theta + \lambda_3 \sin \theta) \tag{8} \\
\mu &= \sqrt{\frac{a_0}{K}} (2\lambda_1 + 2\lambda_2 \sin \theta + 2\lambda_3 \cos \theta) \tag{9} \\
\nu &= \sqrt{\frac{a_0}{K}} (\lambda_4 \sin \theta + \lambda_5 \cos \theta) \tag{10}
\end{align}

Тогда гамильтониан принимает вид:
\begin{equation}
H = l_R \lambda + l_T \mu + l_N \nu
\end{equation}

\subsection{Оптимальное управление}

Согласно принципу максимума Понтрягина, оптимальное управление максимизирует гамильтониан при ограничении:
\begin{equation}
l_R^2 + l_T^2 + l_N^2 = 1
\end{equation}

Решаем задачу на условный экстремум. Функция Лагранжа:
\begin{equation}
\mathcal{L} = l_R \lambda + l_T \mu + l_N \nu + \eta (1 - l_R^2 - l_T^2 - l_N^2)
\end{equation}

Условия стационарности:
\begin{align}
\frac{\partial \mathcal{L}}{\partial l_R} &= \lambda - 2\eta l_R = 0 \\
\frac{\partial \mathcal{L}}{\partial l_T} &= \mu - 2\eta l_T = 0 \\
\frac{\partial \mathcal{L}}{\partial l_N} &= \nu - 2\eta l_N = 0
\end{align}

Отсюда:
\begin{equation}
l_R = \frac{\lambda}{2\eta}, \quad l_T = \frac{\mu}{2\eta}, \quad l_N = \frac{\nu}{2\eta}
\end{equation}

Из условия нормировки:
\begin{equation}
l_R^2 + l_T^2 + l_N^2 = \frac{\lambda^2 + \mu^2 + \nu^2}{4\eta^2} = 1
\end{equation}

Следовательно:
\begin{equation}
2\eta = \sqrt{\lambda^2 + \mu^2 + \nu^2} = p
\end{equation}

Окончательно получаем оптимальное управление:

\begin{align}
l_R &= \frac{\lambda}{p} \tag{7} \\
l_T &= \frac{\mu}{p} \tag{7} \\
l_N &= \frac{\nu}{p} \tag{7}
\end{align}

\section{Преобразование уравнений для праймера}

\subsection{Преобразование радиальной компоненты}

Из уравнения (8):
\begin{equation}
\lambda = \sqrt{\frac{a_0}{K}} (-\lambda_2 \cos \theta + \lambda_3 \sin \theta)
\end{equation}

Введем амплитуду и фазу:
\begin{equation}
A = \sqrt{\lambda_2^2 + \lambda_3^2}, \quad \tan \theta_0 = \frac{\lambda_2}{\lambda_3} \tag{14}
\end{equation}

Тогда:
\begin{align}
\lambda_2 &= A \sin \theta_0 \\
\lambda_3 &= A \cos \theta_0
\end{align}

Подставляем:
\begin{align}
\lambda &= \sqrt{\frac{a_0}{K}} (-A \sin \theta_0 \cos \theta + A \cos \theta_0 \sin \theta) \\
&= \sqrt{\frac{a_0}{K}} A (\cos \theta_0 \sin \theta - \sin \theta_0 \cos \theta) \\
&= \sqrt{\frac{a_0}{K}} A \sin(\theta - \theta_0)
\end{align}

Таким образом:
\begin{equation}
\lambda = \sqrt{\frac{a_0}{K}} \sqrt{\lambda_2^2 + \lambda_3^2} \sin(\theta - \theta_0) \tag{11}
\end{equation}

\subsection{Преобразование трансверсальной компоненты}

Из уравнения (9):
\begin{equation}
\mu = \sqrt{\frac{a_0}{K}} (2\lambda_1 + 2\lambda_2 \sin \theta + 2\lambda_3 \cos \theta)
\end{equation}

Выражаем через введенные переменные:
\begin{align}
\mu &= \sqrt{\frac{a_0}{K}} [2\lambda_1 + 2A(\sin \theta_0 \sin \theta + \cos \theta_0 \cos \theta)] \\
&= \sqrt{\frac{a_0}{K}} [2\lambda_1 + 2A \cos(\theta - \theta_0)]
\end{align}

Таким образом:
\begin{equation}
\mu = \sqrt{\frac{a_0}{K}} \left[2\lambda_1 + 2\sqrt{\lambda_2^2 + \lambda_3^2} \cos(\theta - \theta_0)\right] \tag{12}
\end{equation}

\subsection{Преобразование нормальной компоненты}

Из уравнения (10):
\begin{equation}
\nu = \sqrt{\frac{a_0}{K}} (\lambda_4 \sin \theta + \lambda_5 \cos \theta)
\end{equation}

Это выражение можно представить в виде:
\begin{equation}
\nu = \sqrt{\frac{a_0}{K}} [B \sin(\theta - \theta_0) + C \cos(\theta - \theta_0)]
\end{equation}

где коэффициенты B и C находятся из системы:
\begin{align}
B \cos \theta_0 - C \sin \theta_0 &= \lambda_4 \\
B \sin \theta_0 + C \cos \theta_0 &= \lambda_5
\end{align}

Решаем эту систему:
\begin{align}
B &= \lambda_4 \cos \theta_0 + \lambda_5 \sin \theta_0 \\
C &= -\lambda_4 \sin \theta_0 + \lambda_5 \cos \theta_0
\end{align}

Выражаем через $\lambda_2$ и $\lambda_3$:
\begin{align}
\cos \theta_0 &= \frac{\lambda_3}{\sqrt{\lambda_2^2 + \lambda_3^2}} \\
\sin \theta_0 &= \frac{\lambda_2}{\sqrt{\lambda_2^2 + \lambda_3^2}}
\end{align}

Тогда:
\begin{align}
B &= \frac{\lambda_4 \lambda_3 + \lambda_5 \lambda_2}{\sqrt{\lambda_2^2 + \lambda_3^2}} \\
C &= \frac{-\lambda_4 \lambda_2 + \lambda_5 \lambda_3}{\sqrt{\lambda_2^2 + \lambda_3^2}}
\end{align}

Таким образом:
\begin{equation}
\nu = \sqrt{\frac{a_0}{K}} \left[
\frac{\lambda_3 \lambda_4 - \lambda_2 \lambda_5}{\sqrt{\lambda_2^2 + \lambda_3^2}} \sin(\theta - \theta_0) + 
\frac{\lambda_2 \lambda_4 + \lambda_3 \lambda_5}{\sqrt{\lambda_2^2 + \lambda_3^2}} \cos(\theta - \theta_0)
\right] \tag{13}
\end{equation}

\section{Анализ конфигураций праймера}

\subsection{Общий вид праймера}

Компоненты праймера образуют эллипс в 3D-пространстве. Квадрат модуля праймера:

\begin{align}
p^2 &= \lambda^2 + \mu^2 + \nu^2 \\
&= \frac{a_0}{K} \left\{
\lambda_2^2 + \lambda_3^2) \sin^2(\theta - \theta_0) + 
[2\lambda_1 + 2\sqrt{\lambda_2^2 + \lambda_3^2} \cos(\theta - \theta_0)]^2 + 
\left[\frac{\lambda_3 \lambda_4 - \lambda_2 \lambda_5}{\sqrt{\lambda_2^2 + \lambda_3^2}} \sin(\theta - \theta_0) + \frac{\lambda_2 \lambda_4 + \lambda_3 \lambda_5}{\sqrt{\lambda_2^2 + \lambda_3^2}} \cos(\theta - \theta_0)\right]^2
\right\}
\end{align}

\subsection{Узловой случай (Nodal Case)}

Условие: центр эллипса в начале координат. Это означает, что $\mu$ должно быть нечетной функцией относительно $\theta - \theta_0$, что требует:

\begin{equation}
\lambda_1 = 0 \tag{15}
\end{equation}

При этом праймер принимает вид:
\begin{align}
\lambda &= \sqrt{\frac{a_0}{K}} A \sin(\theta - \theta_0) \\
\mu &= \sqrt{\frac{a_0}{K}} 2A \cos(\theta - \theta_0) \\
\nu &= \sqrt{\frac{a_0}{K}} [B \sin(\theta - \theta_0) + C \cos(\theta - \theta_0)]
\end{align}

Для симметрии максимумов необходимо, чтобы выполнялось:
\begin{equation}
\theta_2 = \theta_1 + \pi \tag{16}
\end{equation}

И компоненты праймера удовлетворяли:
\begin{align}
\lambda(\theta_1) &= -\lambda(\theta_2) \\
\mu(\theta_1) &= -\mu(\theta_2) \\
\nu(\theta_1) &= -\nu(\theta_2) \tag{17}
\end{align}

\subsection{Невырожденный случай (Nondegenerate Case)}

Условия:
\begin{align}
\lambda_2 \lambda_4 &= -\lambda_3 \lambda_5 \tag{18} \\
\theta_2 - \theta_0 &= \theta_0 - \theta_1 \tag{19}
\end{align}

Компоненты праймера удовлетворяют:
\begin{align}
\lambda(\theta_1) &= -\lambda(\theta_2) \\
\mu(\theta_1) &= \mu(\theta_2) \\
\nu(\theta_1) &= -\nu(\theta_2) \tag{20}
\end{align}

Условие для косинуса:
\begin{equation}
\cos(\theta_1 - \theta_0) = \frac{4\lambda_1 (\sqrt{\lambda_2^2 + \lambda_3^2})^3}{(\lambda_3 \lambda_4 - \lambda_2 \lambda_5)^2 - 3(\lambda_2^2 + \lambda_3^2)^2} \tag{21}
\end{equation}

Неравенства:
\begin{align}
3(\lambda_2^2 + \lambda_3^2)^2 &< (\lambda_3 \lambda_4 - \lambda_2 \lambda_5)^2 \tag{22} \\
4\lambda_1 (\sqrt{\lambda_2^2 + \lambda_3^2})^3 &\leq (\lambda_3 \lambda_4 - \lambda_2 \lambda_5)^2 - 3(\lambda_2^2 + \lambda_3^2)^2 \tag{23}
\end{align}

\subsection{Сингулярный случай (Singular Case)}

Условия:
\begin{align}
\lambda_1 &= 0 \tag{24} \\
\lambda_2 \lambda_4 &= -\lambda_3 \lambda_5 \tag{25} \\
3(\lambda_2^2 + \lambda_3^2)^2 &= (\lambda_3 \lambda_4 - \lambda_2 \lambda_5)^2 = \frac{3}{4} \frac{K}{a_0} \tag{26}
\end{align}

Компоненты праймера (при $p = 1$):
\begin{align}
\lambda &= \frac{1}{2} \sin(\theta - \theta_0) \tag{27} \\
\mu &= \cos(\theta - \theta_0) \tag{28} \\
\nu &= \frac{\sqrt{3}}{2} \sin(\theta - \theta_0) \tag{29}
\end{align}

\section{Решение двухточечной краевой задачи}

\subsection{Узловой случай}

Из уравнений (1)-(5) и условий симметрии (17) получаем изменения элементов:

\begin{align}
\Delta x_1 &= \int \frac{dx_1}{du} du = 2\sqrt{\frac{a_0}{K}} \int l_T du \\
&= 2\sqrt{\frac{a_0}{K}} [l_{T_1} u_1 + l_{T_2} u_2]
\end{align}

Но из симметрии: $l_{T_2} = -l_{T_1}$, поэтому:
\begin{equation}
\Delta x_1 = 2\sqrt{\frac{a_0}{K}} l_{T_1} (u_1 - u_2) = 2\sqrt{\frac{a_0}{K}} l_{T_1} u (1 - 2\delta) \tag{32}
\end{equation}
где $\delta = u_2 / u$.

Аналогично для других компонент:
\begin{align}
\Delta x_2 &= -\sqrt{\frac{a_0}{K}} l_{R_1} u \tag{33} \\
\Delta x_3 &= 2\sqrt{\frac{a_0}{K}} l_{T_1} u \tag{34} \\
\Delta x_4 &= 0 \tag{35} \\
\Delta x_5 &= \sqrt{\frac{a_0}{K}} l_{N_1} u \tag{36}
\end{align}

Из (34) находим:
\begin{equation}
l_{T_1} = \frac{\Delta x_3}{2u} \sqrt{\frac{K}{a_0}}
\end{equation}

Из (33):
\begin{equation}
l_{R_1} = -\frac{\Delta x_2}{u} \sqrt{\frac{K}{a_0}}
\end{equation}

Из (36):
\begin{equation}
l_{N_1} = \frac{\Delta x_5}{u} \sqrt{\frac{K}{a_0}}
\end{equation}

Условие нормировки:
\begin{equation}
l_{R_1}^2 + l_{T_1}^2 + l_{N_1}^2 = 1
\end{equation}

Подставляем:
\begin{equation}
\left(\frac{\Delta x_2}{u}\right)^2 \frac{K}{a_0} + \left(\frac{\Delta x_3}{2u}\right)^2 \frac{K}{a_0} + \left(\frac{\Delta x_5}{u}\right)^2 \frac{K}{a_0} = 1
\end{equation}

Умножаем на $u^2$:
\begin{equation}
\frac{K}{a_0} \left(\Delta x_2^2 + \frac{\Delta x_3^2}{4} + \Delta x_5^2\right) = u^2
\end{equation}

Таким образом:
\begin{equation}
u = \sqrt{\frac{K}{a_0}} \sqrt{\Delta i^2 + \frac{\Delta e_x^2}{4} + \Delta e_y^2} \tag{37}
\end{equation}

Величина первого импульса из (32):
\begin{equation}
u_1 = \frac{\Delta e_x + \Delta a/a_0}{2\Delta e_x} \sqrt{\frac{K}{a_0}} \sqrt{\Delta i^2 + \frac{\Delta e_x^2}{4} + \Delta e_y^2} \tag{38}
\end{equation}

\subsection{Невырожденный случай}

Вводим новые переменные:
\begin{align}
R &\equiv \sqrt{\frac{a_0}{K}} \cdot \sqrt{\lambda_2^2 + \lambda_3^2} \tag{50} \\
C &\equiv \cos(\theta - \theta_0) \\
S &\equiv \sin(\theta - \theta_0)
\end{align}

И нормировку:
\begin{equation}
\sqrt{\frac{a_0}{K}} \cdot \frac{\lambda_3 \lambda_4 - \lambda_2 \lambda_5}{\sqrt{\lambda_2^2 + \lambda_3^2}} \equiv 1 \tag{49}
\end{equation}

После преобразований получаем выражения для направляющих косинусов (52)-(54) и изменений элементов (55)-(59).

Далее вводим углы:
\begin{equation}
\tan \omega \equiv \frac{\Delta x_2}{\Delta x_3}, \quad \tan \Omega \equiv \frac{\Delta x_4}{\Delta x_5} \tag{60}
\end{equation}

И переменную:
\begin{equation}
T \equiv \tan(\omega - \Omega) \tag{61}
\end{equation}

После решения системы получаем выражение для суммарного импульса:
\begin{equation}
u = \sqrt{\frac{K}{2a_0}} \sqrt{\Delta i^2 + \Delta e_x^2 + \Delta e_y^2 - \frac{\Delta a^2}{2a_0^2} + \sqrt{\left(\Delta i^2 - \Delta e_x^2 - \Delta e_y^2 + \frac{\Delta a^2}{a_0^2}\right)^2 + 4\Delta i^2 \Delta e_y^2}} \tag{68}
\end{equation}

\subsection{Сингулярный случай}

Для сингулярного случая получаем:
\begin{equation}
u = \sqrt{\frac{K}{4a_0}} \sqrt{\left(\sqrt{3}\Delta i + \Delta e_y\right)^2 + \Delta e_x^2} \tag{75}
\end{equation}

\end{document}